% discussion.tex (국문 v14, 2026-10-20): 소제목 없음(Sci Rep 권장 구성). 여섯 문단: 기제, 제품, 선행 연구, 지도와 새 지역 적용, 한계, 전망.
\section{고찰}\label{sec:discussion}

물리 지식이 모델에 들어가는 방식이, 기계학습이 지역 Stefan 모델을 넘어서는지와 그 개선이 라벨 없는 지역으로 전이되는지를 결정하였다. 물리 잔차 구조(식~(\ref{eq:residual}))에서는 Stefan 항이 각 해 융해 계절 기후에 대한 응답을 한정된 형태로 담고, 학습기는 토양과 지형이 유효 계수에 주는 편차를 보정한다. 네 가지 결과가 이 역할 분담을 뒷받침한다. 이득은 지형 · 토양 공변량을 요구하였고 기후 공변량만으로는 나타나지 않았다. 이득은 Stefan 오차가 큰 곳에서 가장 컸다. 잔차 학습기의 기여도는 레나 델타에서 지형, 캐나다에서 토양 유기탄소를 가리켰다. Stefan 유사라벨은 학습기를 Stefan 모델 쪽으로 끌어당겼을 뿐 정확도를 더하지 않았다. 라벨에서 기후 응답까지 배워야 하는 직접 ML은 한 지역의 라벨을 모두 쓸 때 더 낮은 점 오차에 이르렀으나, 그 지역을 학습에서 빼면 레나 델타 오차가 Stefan 모델보다 \ndrefLeregDzabs{} cm 컸고, 라벨 25개로는 알래스카와 캐나다에서 Stefan 모델보다 부정확하였다. 따라서 잔차 학습기의 이득은 조건부이며 그 조건 자체가 결과이다. 토양 · 지형 정보가 있고 지역 Stefan 계수가 지점들을 잘못 나타내는 곳에서는 물리에 학습기를 더할 가치가 있고, 대상 지역 라벨이 적은 곳에서는 학습기를 물리에 묶어 둘 가치가 있다. 후자는 영구동토 지역의 흔한 상황이다.

같은 셀 · 같은 해의 공개 제품 비교는 제품 정확도가 지역마다 다른 이유를 설명한다. 두 ML 제품은 여러 해에 걸쳐 모은 지점 관측을 다년 예측 변수에 연결하며, 검증은 각자의 지점 집합에서 무작위 또는 지점 단위 분할로 한다\cite{Liu2024,Wei2026}. 학습 지점이 우리 라벨 셀 가까이에 있는 곳(알래스카, 그리고 같은 현장 관측의 이전 판으로 학습한 Liu 등의 경우 캐나다)에서는 제품이 우리 모델만큼 정확하였고, 그곳의 비교는 독립이 아니다. 가까운 학습 지점이 없는 곳(레나 델타, 중앙값 거리 \ndistweiLemed{} km)에서는 제품이 관측 ALT와 무관하게 약 40--70 cm를 부여하였고(상관 \nrLeWei{}, \nrLeLiu{} 대 \nrLeRo), 델타 북부의 얇은 활동층을 약 20 cm 과대 추정하였으며, 레나 델타 라벨 없이 학습한 물리 잔차 모델보다 부정확하였다. 이 차이는 우리 워크플로의 세 가지 성질에서 비롯한다. 각 예측의 기후 입력은 명시된 해의 융해 도일을 관측일까지 누적한 값이므로 추운 해의 얇은 활동층이 두꺼운 해와 평균되지 않는다. Stefan 기준값은 학습 집합 밖의 셀과 해로 기후 응답을 전달한다. 평가는 모든 방법과 제품을 같은 지점-연도에서 제품 학습 지점 마스크와 함께 채점하므로 지역별 순위를 해석할 수 있다. ALT 관측으로 학습하지 않는 과정 기반 제품 ESA CCI는 외부 CALM 지점에서 지역 라벨 5개를 쓴 우리 모델보다 오차가 컸다. 두 접근을 한 척도에 놓고 비교한 것은 이 연구가 처음이다.

이 결과는 선행 비교를 세 방향으로 확장한다. Gautam 등\cite{Gautam2025}은 알래스카 CALM 지점 68곳의 한 해 자료에서 랜덤 포레스트가 홀드아웃 지점에서 정확도를 대부분 잃는 반면 Stefan 모델은 유지함을 보였다. 이 연구는 그 순서를 세 지역의 2010--2024년 지점-연도와 지역 홀드아웃에서 확인하였고, Stefan 모델을 보정하는 학습기가 그 전이성을 유지하므로 두 접근이 양자택일이 아님을 보였다. 호수 · 하천의 물리 기반 학습 연구는 과정 지식이 자료 부족과 외삽 조건에서 주로 도움이 된다고 보고하였다\cite{Read2019,Jia2021,Willard2022}. 라벨 수 실험과 지역 홀드아웃은 이 진술을 ALT에 대해 정량화하여, 이득이 판정되는 라벨 수와 대상 지역 라벨 없는 오차를 제시한다. 끝으로 무작위 교차검증이 직접 ML의 공간 블록 오차를 \noptimLeDz--\noptimCaDz{} cm 과소평가하고 지역 홀드아웃 오차는 그보다 더 크다는 결과는 공간 검증 문헌\cite{Roberts2017,Ploton2020,MeyerPebesma2022}을 따른다. ALT 제품에 대한 함의는, 무작위 · 지점 분할의 오차가 학습 관측이 없는 지역에서의 제품 정확도를 나타내지 못하며, 그런 지역이야말로 제품이 가장 필요한 곳이라는 점이다.

연도별 지도는 이 결과를 사용 조건이 명시된 산출물로 바꾼다. 각 지도는 한 해의 ERA5-Land 기후(9월 1일까지 누적한 TDD)로 구동되므로 같은 해의 관측이나 제품과 비교할 수 있고, 적용 가능 영역과 예측 구간을 지닌다. 지도의 지역 평균 편차는 관측망이 안정된 레나 델타에서 관측 편차를 따랐으나($r$ = \nanobsrLe), 관측 셀 구성이 캠페인마다 바뀌어 관측 편차의 범위가 예측보다 넓은 알래스카에서는 약하게만 따랐다($r$ = \nanobsrAk). 지도의 선형 추세는 2011--2024년 기후에 대한 모델 응답이며 관측 추세가 아니다. 새 지역에 대해서는 라벨 수 실험이 고정된 절차가 아니라 기대할 수 있는 바를 알려 준다. 지역 라벨이 없으면 다른 지역의 계수를 쓴 Stefan 모델이나 다른 지역으로 학습한 물리 잔차 모델이 더 안전한 출발점이다. 지역 라벨은 그 지역의 토양 · 지형 범위를 대표할 때에만 모델을 개선한다. 레나 델타 라벨 25개는 물리 잔차 모델의 오차를 \ntlrmseLeRszero{} cm에서 \ntlrmseLeRstwfive{} cm로, 외부 CALM 라벨 5개는 \nftAllRszero{} cm에서 \nftAllRsfive{} cm로 줄였지만, 캐나다 라벨 25개는 캐나다 지점의 융해 깊이 범위가 넓어 재보정 Stefan 모델의 오차를 \ntlrmseCaPzzero{} cm에서 \ntlrmseCaPntwfive{} cm로 키웠다. 새 지역 안에서의 공간 블록 평가와 결과 지도의 적용 가능 영역이 이 두 경우를 가르는 점검이며, 둘 다 워크플로에 포함되어 있다.

이 결론에는 한계가 있다. 알래스카 라벨의 절반이 2014년 자료여서 이 지역에서는 관측 연도 입력의 효과를 판정하지 못하였고, 레나 델타 라벨은 블록 \nnblkLe{}개의 자료 하나에서 와서 신뢰 구간이 넓다. 라벨은 독립된 1 km 셀이 아니라 관측 지점이다. 한 횡단선의 여러 지점이 한 셀에 들어갈 수 있으며, 공간 블록 설계가 이 의존성을 검증에서 제거하더라도 지점-연도 수는 지도 규모의 독립 표본 수를 과대 표시한다. 기후 입력은 월 단위이고 토양 공변량은 정적이어서 단기 기상과 토양 수분 변화는 표현되지 않는다. 주로 산지와 영구동토 지대 남쪽 가장자리의 약 3 m 넘는 활동층은 학습 범위 밖이며, 외부 시험은 지역 라벨 없이는 잔차 학습기가 그런 조건으로 전이되지 않음을 보였다. 공통 연도 한 해의 지도 · 관측 비교는 레나 델타와 캐나다에서 적은 셀에 근거한다.

일 단위 기후 입력, 원격탐사 적설 · 토양 수분 정보, CALM과 캠페인 자료가 함께 있는 서시베리아 같은 네 번째 지역의 라벨이 다음 단계이다. 이 연구의 모든 수치는 명시된 지점 · 연도 · 검증 설계에 묶여 있으므로, 같은 프로토콜을 새로 나오는 제품과 지역에 적용할 수 있고, 연도별 지도는 기후 자료가 공개되는 대로 해마다 이어 붙일 수 있다.
